\documentclass[11pt]{article}

\usepackage[margin=0.85in, top=0.7in, bottom=0.7in]{geometry}
\usepackage{graphicx}
\usepackage{microtype}
\usepackage{parskip}
\usepackage{caption}
\usepackage{hyperref}

\hypersetup{colorlinks=true, linkcolor=blue, citecolor=blue}

\setlength{\parskip}{4pt}

\title{\vspace{-1em}\textbf{Milestone 3: ARIA -- Adaptive Runtime Index with Async-flush}\vspace{-0.5em}}
\author{Md Toki Tahmid \quad COS568, Princeton University}
\date{}

\begin{document}
\maketitle
\vspace{-1em}

\section*{Approach}

The Milestone 2 naive hybrid flushes DPGM into LIPP synchronously, stalling
every insert that triggers the threshold for the full duration of the migration.
ARIA fixes this with two ideas: (1) the flush runs on a background thread so
inserts are never blocked, and (2) the index decides at runtime whether to use
the write buffer at all.

\textbf{Runtime profiling.} During the first 1024 operations after bulk load,
ARIA counts inserts and lookups. If inserts exceed 50\%, it stays in
\texttt{WRITE\_BUFFERED} mode; otherwise it drains the DPGM buffer into LIPP
once and switches to \texttt{READ\_OPTIMIZED} mode permanently. The profiling
flag is cleared after this decision and never checked again.

\textbf{WRITE\_BUFFERED mode.} Inserts go into \texttt{dpgm\_active\_}.
When it reaches \texttt{flush\_interval} keys (2M), ARIA atomically moves it
to \texttt{dpgm\_flushing\_}, resets \texttt{dpgm\_active\_} to an empty
buffer, and spawns a background thread to drain \texttt{dpgm\_flushing\_} into
LIPP one key at a time. Incoming inserts continue into the fresh buffer without
waiting. Lookups check LIPP, then \texttt{dpgm\_active\_}, then
\texttt{dpgm\_flushing\_}. A \texttt{shared\_mutex} protects LIPP during the
flush: the background thread holds a \texttt{unique\_lock}; concurrent lookups
acquire a \texttt{shared\_lock} only when the \texttt{flush\_in\_progress\_}
atomic flag (with acquire/release ordering) is set.

\textbf{READ\_OPTIMIZED mode.} Both inserts and lookups go directly to LIPP.
The hot path is a single predicted branch with no locking overhead, equivalent
to running pure LIPP. The internal LIPP instance (MyLIPP) adds one
\texttt{\_\_builtin\_expect} hint to the \texttt{find()} traversal loop,
giving roughly 2.6\% higher throughput over unmodified LIPP on the FB dataset.

\textbf{No keys lost.} Before switching to READ\_OPTIMIZED, all keys in
\texttt{dpgm\_active\_} are inserted into LIPP. During an async flush, keys
in \texttt{dpgm\_flushing\_} remain reachable via the lookup path until the
thread finishes. No key is ever discarded.

\textbf{Hyperparameters.} \texttt{pgm\_error} was swept over \{64, 128, 256\}
with \texttt{flush\_interval} fixed at 2M. The variant with the highest average
throughput across 3 runs was selected per workload.

\section*{Results}

Figures~\ref{fig:fb}--\ref{fig:osmc} show throughput (Mops/s) and index size
(GB) for both workloads across all three datasets. ARIA (labeled
\textit{HybridPGMLIPPAsync}) outperforms DPGM, LIPP, and the M2 naive hybrid
on both workloads across all datasets.

\begin{figure}[h]
    \centering
    \includegraphics[width=0.82\linewidth]{../analysis_results/milestone3_fb.png}
    \caption{Facebook (FB) dataset.}
    \label{fig:fb}
\end{figure}

\begin{figure}[h]
    \centering
    \includegraphics[width=0.82\linewidth]{../analysis_results/milestone3_books.png}
    \caption{Books dataset.}
    \label{fig:books}
\end{figure}

\begin{figure}[h]
    \centering
    \includegraphics[width=0.82\linewidth]{../analysis_results/milestone3_osmc.png}
    \caption{OSMC dataset.}
    \label{fig:osmc}
\end{figure}

\end{document}
